\chapter{Wound Debridement}

\section*{Overview}
Wound debridement is the medical removal of dead, damaged, or infected tissue to improve the healing potential of the remaining healthy tissue.

\section*{Alternative Names}
Surgical debridement, Necrotic tissue removal

\section*{Principle}
Dead, infected, or foreign tissue is removed from a wound by surgical cutting, enzymes, or other means, leaving a clean viable bed. Removing necrotic tissue eliminates the medium for bacterial growth and allows granulation and healing to proceed.
\section*{History}
Developed in military surgery. Ambroise Pare in the 16th century advocated against cauterizing wounds and for cleaning devitalized tissue.

\section*{Why This Test or Procedure Is Ordered}
Ordered for chronic, non-healing wounds, diabetic foot ulcers, pressure sores, or heavily contaminated traumatic wounds.

\section*{Is This a Primary or Secondary Test or Procedure}
Primary treatment for chronic wounds and necrotizing soft-tissue infections.

\section*{How This Test or Procedure Is Performed}
Can be performed surgically (using scalpel/scissors to cut away dead tissue), enzymatically (using chemical ointments), or autolytically (using special dressings).

\section*{What Preparation Is Needed}
Administer local or systemic pain medication. Irrigate the wound.

\section*{Contraindications and Precautions}
Severe arterial insufficiency (where dry gangrene should be left intact until revascularization is achieved).

\section*{Risks}
Pain, bleeding, damage to underlying structures (tendons, nerves), or introducing infection if sterile technique is breached.

\section*{What Happens After the Test or Procedure}
Apply a sterile dressing designed for wound healing (alginate, hydrogel). Monitor wound dimensions and drainage.

\section*{Understanding Your Results}
Success is marked by a clean, red wound bed composed of healthy granulation tissue with no slough or eschar.

\section*{Normal Results}
Clean wound bed containing healthy red granulation tissue; no purulent odor or necrotic tissue.

\section*{Follow-up Tests and Procedures}
Frequent dressing changes, regular wound measurements, and weekly debridement sessions as needed.

\section*{References}
\begin{enumerate}
  \item MedlinePlus. Wound Debridement. U.S. National Library of Medicine; 2025.
  \item Current Medical Diagnosis and Treatment. McGraw-Hill; 2025.
\end{enumerate}

\clearpage
